\section{Area overlap: a companion-kernel check against the mock's own
algorithm}
\label{sec:area_overlap}

Section~\ref{sec:results}'s $\Delta^{\rm prj}_{\rm RND} = P_1 + \beff I_2$
agreement (\DeltaPrjAgreeMin{}--\DeltaPrjAgreeMax{}) still leaves a
per-bin residual too large to attribute to the MOR fix alone (issue \#5).
This section isolates one candidate: `area\_overlap`, the closed-form
disk-overlap fraction weighting every companion halo's contribution to
$P_1$, $I_1$, $I_2$.

\pyfile{clenspy.selection.geometry.area\_overlap($\theta$, $\theta_{\lob}$,
$\theta_{\ltr}$)} treats a companion halo's richness as smeared
\emph{uniformly} over its own aperture disk of angular radius
$\theta_{\ltr}$, and returns the geometric overlap area of that disk with
the target's disk of radius $\theta_{\lob}$. That is not what the mock
actually does. Reading its generator verbatim
(\pyfile{SelectionBias/make\_mock\_lob\_sigma\_catalog.py:run\_lambda\_ob},
faithful to notebook cell 24): every candidate cluster's
$\lambda^{\rm ob}$ is measured by (i) a hard angular cut, keeping only
galaxies with $\theta < \theta_{\lob}$, and (ii) among those, a soft
radial weight $\theta_r(x) = \tfrac12\,{\rm erfc}\big(15(x-1.2)\big)$,
$x = R_i/R_{\lambda,{\rm target}}$ -- but the hard cut already confines
$x < 1$, where $\theta_r \ge 0.9999$, i.e.\ effectively $1$. So the real
selection is purely geometric: what fraction of a companion's own
galaxies -- one central at its centre plus
$\mathrm{round}(\ltr) - 1$ satellites drawn from its own projected-NFW
profile (the same \pyfile{Mprj\_com}/\pyfile{fun} kernel as the mock's
$\Sigma_{\rm NFW}$, notebook cell 15) -- fall inside the \emph{target's}
aperture disk.

\subsection{An exact, unit-free reduction}

At fixed redshift $\theta = R_\lambda(1+z)/\chi(z)$ for both objects, so
every angle ratio equals the corresponding $R_\lambda$ ratio and the
$(1+z)/\chi(z)$ factors cancel identically between target and companion.
Working in target-aperture units ($\theta_{\lob} \equiv 1$), the
companion's own aperture is $r_{\ltr} = \theta_{\ltr}/\theta_{\lob}$ and
its angular separation from the target is $s = \theta_{\rm sep}/\theta_{\lob}$
-- both dimensionless, both independent of $h$, $z$, and $\chi(z)$. The
true overlap fraction is then a deterministic ring integral, not a Monte
Carlo estimate: with the companion's satellites radially distributed as
$p(r)\,dr \propto r\,\Sigma_{\rm NFW}(r)\,dr$ on $r \in [0, r_{\ltr}]$
(concentration $c$, scale radius $r_{\ltr}/c$), and
${\rm ring}(r, s)$ the exact angular fraction of a ring of radius $r$
centred a distance $s$ from a unit-radius target disk that lies inside
it,
\[
f_{\rm true}(s, r_{\ltr}, c) = \frac{
  \mathbb{1}[s<1] + (\ltr - 1)\int_0^{r_{\ltr}} p(r)\,{\rm ring}(r,s)\,dr
}{\ltr},
\qquad
\ltr = \lambda^{\rm ob}\,r_{\ltr}^{5}
\]
(the last identity from $R_\lambda \propto \lambda^{0.2}$). The central
term uses $\lambda^{\rm ob} = 24$, this report's reference-bin richness;
the mock's own generator always adds a deterministic $+1$ central
regardless of mass, so the split is only meaningful once
$\ltr \gtrsim 2$ -- which is why the comparison below starts at
$r_{\ltr} = 0.6$.

\begin{figure}[t]
  \centering
  \includegraphics[width=0.72\textwidth]{fig_area_overlap.pdf}
  \caption{The true companion-richness fraction landing in the target
  aperture, $f_{\rm true}(s)$ (solid, exact NFW ring quadrature,
  $c=\AreaOverlapCFidLow[]$--$\AreaOverlapCFidHigh[]$ tested,
  $c=5$ shown), against \modname{area\_overlap}'s uniform-disk
  approximation (dashed), for three companion sizes relative to the
  target's own aperture.}
  \label{fig:area_overlap}
\end{figure}

\subsection{Result}

Figure~\ref{fig:area_overlap} shows a clean, physically sensible
crossover at $s = 1$: for $s < 1$ (companion centred inside the target's
aperture -- the region that dominates the $\sin\theta$-weighted $P_1$,
$I_1$, $I_2$ integrals), \modname{area\_overlap} \emph{under}-predicts
the true, NFW-concentrated overlap by up to \AreaOverlapMaxUnder{}: a
companion's richness is centrally concentrated at its own core, so more
of it lands in the target's aperture than a uniform smear would give.
For $s > 1$ (companion centred outside), the sign flips and
\modname{area\_overlap} \emph{over}-predicts by up to
\AreaOverlapMaxOver{}: the uniform disk still has mass reaching toward
the target from its near edge, where the NFW profile has already fallen
off. The mean absolute deviation across the full $s$ range tested is
\AreaOverlapMeanAbs{}. This is not a concentration-tuning artefact --
varying $c$ from \AreaOverlapCFidLow{} to \AreaOverlapCFidHigh{} at
$r_{\ltr}=0.75$ moves the peak deviation only from
\AreaOverlapConcMin{} to \AreaOverlapConcMax{}, i.e.\ the sign and rough
size of the effect are set by the uniform-disk-versus-NFW assumption
itself, not by the concentration value chosen.

Because $P_1$, $I_1$, $I_2$ integrate this kernel against $\sin\theta$
over the \emph{entire} $\theta$ range from the exclusion boundary to
$2\theta_{\lob}$ -- a region straddling $s=1$ -- the under- and
over-prediction partially cancel rather than adding in one direction,
which is consistent with $P_1$ itself checking out to sub-percent
accuracy against an independent closed-form full-range integral
(Sec.~\ref{sec:results}) while the bin-by-bin
$\Delta^{\rm prj}_{\rm RND}/\Delta^{\rm prj}_{\rm mock}$ residual remains
at the \DeltaPrjAgreeMax{} level. The net effect on $I_1$/$I_2$, which
carry an additional $\xi_{\rm NL}(\theta)$ weight that is not flat across
$\theta$, is not yet isolated from this check alone and is the natural
next step: reweighting the ring integral by $\xi_{\rm NL}$ the way the
production operators do, rather than treating $P_1$'s flat weight as
representative of all three operators.

All numbers and the figure in this section are generated by
\pyfile{scripts/make\_area\_overlap\_check.py} $\to$
\pyfile{scripts/make\_figures.py} / \pyfile{scripts/make\_values.py},
independently of \pyfile{scripts/make\_results.py} -- this is a check of
one geometric kernel in isolation, not a rerun of the full pipeline.
